\section{Discussion}
\label{sec:discussion}

\subsection{What the Protocol Contributes}

The push/pause/hold protocol is the contribution---not the CFAR
detector. CFAR is a well-established technique from radar signal
processing~\citep{finn1968cfar}. What is new is the choreography:
autonomous agents, each with their own objectives and unaware of
each other, take turns as active perturbers and passive observers.
The coordination is minimal---a scheduling constraint (one sender at
a time), not information exchange. Agents do not share objectives,
results, or detected edges.

Without the protocol, there is no signal for CFAR to detect. Two
independently adjusting agents produce confounded objective
trajectories that passive methods cannot separate~\citep{sugihara2012ccm}.
The protocol converts the live system into a sequence of structured
experiments.

\subsection{Detection Boundaries Are Parameterized}

The detection boundaries presented in
Tables~\ref{tab:strength}--\ref{tab:noise} are specific to the
observation parameters listed in Section~\ref{sec:obs-params}. The
method has no fundamental sensitivity limit---the boundary is a
function of observation budget (block sizes, number of rounds) and
detector configuration ($P_{\text{fa}}$, minimum sample counts).
The CFAR scale factor $k$ decreases with more reference cells;
the MAD estimate stabilizes with more samples.

We present results at one operating point ($B = 15$, $P_{\text{fa}} =
0.05$) with one budget contrast ($B = 30$ at noise $0.10$). This
characterizes where the method works and where it stops under
reasonable observation costs. The noise failure at $\sigma = 0.10$
represents an SNR of approximately 2.5:1---an aggressive operating
point for live systems.

\subsection{Threshold Calibration}

The CFAR threshold (\ref{eq:threshold}) is computed from per-sample
noise (MAD of individual objective values), but the detection test
compares block medians. The actual false alarm rate is therefore
lower than the nominal $P_{\text{fa}} = 0.05$---the threshold is
conservative. This is confirmed empirically: zero false positives
were observed across all noise levels (Table~\ref{tab:fpr}). A
median-aware threshold would improve sensitivity at high noise but
requires re-deriving the false alarm rate for median statistics.
We leave this to future work and report results with the conservative
threshold.

\subsection{Nonlinear Coupling}

The protocol is coupling-shape-agnostic by design. The sender pushes
to extreme parameters and returns to neutral; the detector asks
whether the receiver's signal changed and recovered. The shape of
the coupling function determines the step amplitude but not whether
a step exists. All four tested shapes (linear, threshold, saturation,
polynomial) were detected at adequate coupling strength
(Table~\ref{tab:nonlinear}).

At low coupling ($0.3$), nonlinear shapes that compress the step
(saturation, polynomial) fall below the detection floor, while
threshold---which produces a large discontinuous jump---is detected.
This is not a limitation of the protocol but a property of the
specific coupling function's step response at the operating point.

\subsection{Topology Limitation}

The scanner discovers step-change coupling between node pairs. It
does not distinguish direct edges from transitive paths. In a
branching topology (e.g., $A \to B \to D$ and $A \to C \to D$),
node $D$ responds to $A$'s perturbation through transitive paths,
and the scanner would report a step change for the pair $(A, D)$
even though no direct edge exists. Distinguishing direct from
transitive coupling requires either known edge weights to subtract
expected contributions or graph-theoretic analysis on the discovered
step-change graph. This is an open problem for future work.

\subsection{What This Is Not}

This is not causal discovery in the Pearl sense~\citep{pearl2009causality}.
Pearl's framework discovers Boolean causal arrows from interventions
designed by a central researcher. Our agents are not researchers;
they are agents with their own objectives who sacrifice
exploration time to cooperate for measurement. The coordination
is scheduling, not experimental design.

This is also not system identification~\citep{ljung1999sid}, which
requires a specified model structure and estimates parameters within
it. Our method discovers whether coupling exists and measures its
step amplitude, without assuming a model.
